\section{Discussão}\label{sec:discussao}
A demonstração evidencia separação entre interação, composição de dados e
persistência. Repositórios e injeção permitiram controlar falhas nos testes;
isso não demonstra redução futura de defeitos. A presença de ViewModels
também não impede acúmulo de responsabilidades. O colaborador de estado evita
repetição, ao custo de mais uma abstração de apresentação.

Quatro módulos preservam isolamento de funcionalidades com menos configuração
que o original. Em contrapartida, as camadas dentro do núcleo têm limites de
pacotes, mais fracos que módulos de compilação independentes. A escolha atende
ao escopo acadêmico, sem pretender substituir a modularização do NiA.

Os resultados são exploratórios e descrevem um emulador em uma máquina,
com APK debug, cache de disco preservado e poucas amostras. Não houve comparação controlada de
arquiteturas nem estudo com usuários. Rótulos e áreas de toque verificáveis
não comprovam usabilidade global. As hipóteses sobre stakeholders também
precisariam de validação com participantes reais.
